\documentclass[a4paper]{article}
\usepackage{amsfonts}
\usepackage{amsmath}
\usepackage{amssymb}
\usepackage{graphicx}     

\begin{document}
  
\noindent \large Auteur: Sabawoon Enayat \\
\noindent \large Studierichting: Informatica\\
\noindent \large Oefeningenreeks 1C nr. 48.1\\

\newcommand{\Bgtan}{\operatorname{Bgtan}}
\newcommand{\cis}{\operatorname{cis}}

\medskip

\normalsize

Opgave: Bereken $(1 + i)^n - (1 - i)^n$, met $n = 2025$\\

$(1 + i)^{2025} - (1 - i)^{2025}$\\

$z_1 = 1 + i$\\

$\Rightarrow r = \sqrt{a^2 + b^2} = \sqrt{1^2 + 1^2} = \sqrt{2}$\\

$\Rightarrow \tan \theta = \dfrac{b}{a} = \dfrac{1}{1} = 1 \Rightarrow \theta = \Bgtan (1) = \dfrac{\pi}{4}$\\

$\Rightarrow z_1 = r \cis \theta = \sqrt{2} \cis \left(\dfrac{\pi}{4}\right)$\\

$z_2 = 1 - i$\\

$\Rightarrow r = \sqrt{a^2 + b^2} = \sqrt{1^2 + (-1)^2} = \sqrt{2}$\\

$\Rightarrow \tan \theta = \dfrac{b}{a} = -\dfrac{1}{1} = -1 \Rightarrow \theta = \Bgtan (-1) = -\dfrac{\pi}{4}$\\

$\Rightarrow z_2 = r \cis \theta = \sqrt{2} \cis \left(-\dfrac{\pi}{4}\right)$\\

$\overset{\text{De Moivre}}{\Rightarrow} (r \cis \theta)^n = r^n \cis (n\theta)$\\

$\left(\sqrt{2}^{2025} \cis \left(2025 \dfrac{\pi}{4}\right)\right) - \left(\sqrt{2}^{2025} \cis \left(-2025 \dfrac{\pi}{4}\right)\right)$\\

$= \sqrt{2}^{2025} \left(\cis\left(2025 \dfrac{\pi}{4}\right) - \cis \left(-2025 \dfrac{\pi}{4}\right)\right)$\\

$= \sqrt{2^{2024} \cdot 2} \left(\cis\left(2024 \dfrac{\pi}{4} + \dfrac{\pi}{4}\right) - \cis \left(-2024 \dfrac{\pi}{4} - \dfrac{\pi}{4}\right)\right)$\\

$= 2^{1012} \cdot \sqrt{2} \left(\cis \left(\dfrac{\pi}{4}\right) - \cis \left(- \dfrac{\pi}{4}\right)\right)$\\

$= 2^{1012} \cdot \sqrt{2} \left(\cos \left(\dfrac{\pi}{4}\right) + i \cdot \sin \left(\dfrac{\pi}{4}\right) - \cos \left(- \dfrac{\pi}{4}\right) - i \cdot \sin \left(- \dfrac{\pi}{4}\right)\right)$\\

$= 2^{1012} \cdot \sqrt{2} \left(\dfrac{\sqrt{2}}{2} + i \cdot \dfrac{\sqrt{2}}{2} - \dfrac{\sqrt{2}}{2} + i \cdot \dfrac{\sqrt{2}}{2}\right)$\\

$= 2^{1012} \cdot \sqrt{2} \cdot (i\sqrt{2}) = i \cdot 2^{1012} \cdot \sqrt{2 \cdot 2} = i \cdot 2^{1012} \cdot 2$\\

$= 2^{1013} \cdot i$\\

\begin{thebibliography}{999}
%\bibitem{a} Referentie
\end{thebibliography}
\end{document} 
